\scanpage{177}\mspage{32}
où ``pris à la lettre'' $p''$ est maintenant une inclusion, mais ``moralement'' c'est \inserted{\unsure{tjrs|tjrs.}} une fibration à fibres très \unsure{polies|jolies} [de dim finie notamment, voire des variétés non singulières - c'est donc beaucoup mieux que ce que nous donne le yoga de Cartan-Serre ``toute application continue est équivalente à une fibration''\ldots] C'est ce dernier diagramme qui est susceptible d'une construction directe, purement algébrique, dans le cadre des schémas, lorsque l'on a développé la construction du voisinage tubulaire étale.

Le point où je veux en venir, c'est que la considération \inserted{de la \unsure{cat.|cat}} des faisceaux (ou \inserted{complexes d'iceux, ou} $n$-champs, \unsure{\ldots|..}) $\unsure{\mathcal P|P}$-constructibles sur un $X$, où $\mathcal P$ est une stratification ``équisingulière'' donnée, revient dans le cas particulier actuel à la connaissance d'un diagramme de \uline{types d'homotopie} \inserted{ordinaires} \inserted{(ou de pro-types, si on revient à la top. étale)}
\diagram{\begin{tikzcd}[ampersand replacement=\&,column sep=large,row sep=large]
\unsure{\Delta|\triangle} \arrow[r,"\unsure{j|J}"] \arrow[d,"p"'] \& X_1\\
X_0 \&
\end{tikzcd}}
\unsure{en|et} prenant des systèmes de \unsure{coeff|coeff.} locaux (ou de $n$-champs loc. constants) sur ces sommets \inserted{$X_0,X_1$ \deleted{$\unsure{\Delta|X}$}} reliés par des hom de compatibilité du type $p^*(F_0)\to j^*\unsure{(}F_1)$. Ce serait un exercice amusant (que je n'\unsure{ai} pas encore fait) de
% archive sheet number: 177
